% =====================================================================
%  Anatomía de un LLM — Fase I — Módulo 5
%  Probabilidad y estadística
%  Compilar: latexmk -r ../latexmkrc M05-probabilidad.tex
% =====================================================================
\documentclass[11pt,a4paper,twoside,openany]{book}
\usepackage{../estilo-apuntes}

\newcommand{\ApuntesCurso}{Anatomía de un LLM}
\newcommand{\ApuntesModulo}{Módulo 5}

\usetikzlibrary{decorations.pathreplacing}

% Anclas de hiperenlace únicas por tema para los enunciados
\makeatletter
\@for\nombre:=definicion,teorema,proposicion,lema,corolario,ejemplo,observacion\do{%
  \expandafter\xdef\csname theH\nombre\endcsname{\noexpand\thechapter.\noexpand\arabic{definicion}}}
\makeatother

% Notación propia del módulo
\DeclareMathOperator{\Var}{Var}
\DeclareMathOperator{\Cov}{Cov}
\DeclareMathOperator{\Ber}{Ber}
\DeclareMathOperator{\Bin}{Bin}
\DeclareMathOperator{\Cat}{Cat}
\DeclareMathOperator{\Unif}{Unif}
\DeclareMathOperator{\Expo}{Exp}
\DeclareMathOperator{\BetaD}{Beta}
\DeclareMathOperator{\Dir}{Dir}
\DeclareMathOperator*{\argmax}{arg\,max}
\DeclareMathOperator*{\argmin}{arg\,min}
\newcommand{\N}{\mathcal{N}}
\newcommand{\PP}{\mathbb{P}}
\newcommand{\iid}{i.i.d.}
\newcommand{\Xbar}{\overline{X}}

\hypersetup{pdftitle={Módulo 5. Probabilidad y estadística (resumen)},pdfauthor={Anatomía de un LLM}}

\makeatletter
\newlabel{ej:spam}{{1.8}{2}{Filtro de correo no deseado}{ejemplo.1.8}{}}
\newlabel{ej:spam@cref}{{[ejemplo][8][]1.8}{[1][2][]2}}
\newlabel{fig:arbol}{{1.1}{3}{Árbol de probabilidades del \cref {ej:spam}. La evidencia es $\PP (G)=0{,}10+0{,}04=0{,}14$ y la probabilidad a posteriori, $0{,}10/0{,}14\approx 0{,}714$}{figure.1.1}{}}
\newlabel{fig:arbol@cref}{{[figure][1][1]1.1}{[1][3][]3}}
\newlabel{eje:1-inversa}{{1.1}{4}{}{ejercicio.1.1}{}}
\newlabel{eje:1-inversa@cref}{{[ejercicio][1][1]1.1}{[1][4][]4}}
\newlabel{eje:1-disjuntos}{{1.2}{4}{}{ejercicio.1.2}{}}
\newlabel{eje:1-disjuntos@cref}{{[ejercicio][2][1]1.2}{[1][4][]4}}
\newlabel{eje:1-detector}{{1.3}{4}{}{ejercicio.1.3}{}}
\newlabel{eje:1-detector@cref}{{[ejercicio][3][1]1.3}{[1][4][]4}}
\newlabel{eje:1-cribado}{{1.4}{4}{}{ejercicio.1.4}{}}
\newlabel{eje:1-cribado@cref}{{[ejercicio][4][1]1.4}{[1][4][]4}}
\newlabel{eje:1-monedas}{{1.5}{4}{}{ejercicio.1.5}{}}
\newlabel{eje:1-monedas@cref}{{[ejercicio][5][1]1.5}{[1][4][]4}}
\newlabel{eje:1-complementos}{{1.6}{5}{}{ejercicio.1.6}{}}
\newlabel{eje:1-complementos@cref}{{[ejercicio][6][1]1.6}{[1][4][]5}}
\newlabel{ej:dado}{{2.7}{7}{Dado equilibrado}{ejemplo.2.7}{}}
\newlabel{ej:dado@cref}{{[ejemplo][7][]2.7}{[1][7][]7}}
\newlabel{ej:tresvalores}{{2.10}{7}{}{ejemplo.2.10}{}}
\newlabel{ej:tresvalores@cref}{{[ejemplo][10][]2.10}{[1][7][]7}}
\newlabel{fig:inversa}{{2.1}{9}{Muestreo por transformada inversa del \cref {ej:muestreo}: el intervalo $[0,1]$ se divide en tramos de longitud $p_k$ y se devuelve el tramo en el que cae $u$}{figure.2.1}{}}
\newlabel{fig:inversa@cref}{{[figure][1][2]2.1}{[1][9][]9}}
\newlabel{ej:muestreo}{{2.16}{9}{Muestreo de un token}{ejemplo.2.16}{}}
\newlabel{ej:muestreo@cref}{{[ejemplo][16][]2.16}{[1][9][]9}}
\newlabel{eje:2-categorica}{{2.1}{9}{}{ejercicio.2.1}{}}
\newlabel{eje:2-categorica@cref}{{[ejercicio][1][2]2.1}{[1][9][]9}}
\newlabel{eje:2-dados}{{2.2}{9}{}{ejercicio.2.2}{}}
\newlabel{eje:2-dados@cref}{{[ejercicio][2][2]2.2}{[1][9][]9}}
\newlabel{eje:2-softmax}{{2.3}{10}{}{ejercicio.2.3}{}}
\newlabel{eje:2-softmax@cref}{{[ejercicio][3][2]2.3}{[1][9][]10}}
\newlabel{eje:2-dados4}{{2.4}{10}{}{ejercicio.2.4}{}}
\newlabel{eje:2-dados4@cref}{{[ejercicio][4][2]2.4}{[1][10][]10}}
\newlabel{eje:2-binomial}{{2.5}{10}{}{ejercicio.2.5}{}}
\newlabel{eje:2-binomial@cref}{{[ejercicio][5][2]2.5}{[1][10][]10}}
\newlabel{eje:2-covarianza}{{2.6}{10}{}{ejercicio.2.6}{}}
\newlabel{eje:2-covarianza@cref}{{[ejercicio][6][2]2.6}{[1][10][]10}}
\newlabel{fig:normal}{{3.1}{13}{Densidad de la normal estándar. La zona oscura, $\abs {z}\le 1$, contiene el $68{,}3\,\%$ de la probabilidad; junto con las zonas claras, $\abs {z}\le 2$, el $95{,}4\,\%$}{figure.3.1}{}}
\newlabel{fig:normal@cref}{{[figure][1][3]3.1}{[1][12][]13}}
\newlabel{ej:normal}{{3.9}{13}{}{ejemplo.3.9}{}}
\newlabel{ej:normal@cref}{{[ejemplo][9][]3.9}{[1][12][]13}}
\newlabel{eje:3-tabla}{{3.1}{14}{}{ejercicio.3.1}{}}
\newlabel{eje:3-tabla@cref}{{[ejercicio][1][3]3.1}{[1][14][]14}}
\newlabel{eje:3-espera}{{3.2}{14}{}{ejercicio.3.2}{}}
\newlabel{eje:3-espera@cref}{{[ejercicio][2][3]3.2}{[1][14][]14}}
\newlabel{eje:3-cuadrado}{{3.3}{14}{}{ejercicio.3.3}{}}
\newlabel{eje:3-cuadrado@cref}{{[ejercicio][3][3]3.3}{[1][14][]14}}
\newlabel{eje:3-init}{{3.4}{15}{}{ejercicio.3.4}{}}
\newlabel{eje:3-init@cref}{{[ejercicio][4][3]3.4}{[1][14][]15}}
\newlabel{eje:3-exponencial}{{3.5}{15}{}{ejercicio.3.5}{}}
\newlabel{eje:3-exponencial@cref}{{[ejercicio][5][3]3.5}{[1][15][]15}}
\newlabel{eje:3-he}{{3.6}{15}{}{ejercicio.3.6}{}}
\newlabel{eje:3-he@cref}{{[ejercicio][6][3]3.6}{[1][15][]15}}
\newlabel{ej:emv-bernoulli}{{4.5}{17}{Bernoulli}{ejemplo.4.5}{}}
\newlabel{ej:emv-bernoulli@cref}{{[ejemplo][5][]4.5}{[1][16][]17}}
\newlabel{ej:emv-normal}{{4.7}{17}{}{ejemplo.4.7}{}}
\newlabel{ej:emv-normal@cref}{{[ejemplo][7][]4.7}{[1][17][]17}}
\newlabel{eje:4-ceros}{{4.1}{19}{}{ejercicio.4.1}{}}
\newlabel{eje:4-ceros@cref}{{[ejercicio][1][4]4.1}{[1][19][]19}}
\newlabel{eje:4-datos}{{4.2}{19}{}{ejercicio.4.2}{}}
\newlabel{eje:4-datos@cref}{{[ejercicio][2][4]4.2}{[1][19][]19}}
\newlabel{eje:4-moneda}{{4.3}{19}{}{ejercicio.4.3}{}}
\newlabel{eje:4-moneda@cref}{{[ejercicio][3][4]4.3}{[1][19][]19}}
\newlabel{eje:4-u2}{{4.4}{19}{}{ejercicio.4.4}{}}
\newlabel{eje:4-u2@cref}{{[ejercicio][4][4]4.4}{[1][19][]19}}
\newlabel{eje:4-u3}{{4.5}{19}{}{ejercicio.4.5}{}}
\newlabel{eje:4-u3@cref}{{[ejercicio][5][4]4.5}{[1][19][]19}}
\newlabel{eje:4-sesgo}{{4.6}{20}{}{ejercicio.4.6}{}}
\newlabel{eje:4-sesgo@cref}{{[ejercicio][6][4]4.6}{[1][19][]20}}
\newlabel{eje:4-beta}{{4.7}{20}{}{ejercicio.4.7}{}}
\newlabel{eje:4-beta@cref}{{[ejercicio][7][4]4.7}{[1][20][]20}}
\newlabel{ej:pi}{{5.6}{22}{Estimación de $\pi $}{ejemplo.5.6}{}}
\newlabel{ej:pi@cref}{{[ejemplo][6][]5.6}{[1][22][]22}}
\newlabel{ej:benchmark}{{5.11}{23}{Banco de pruebas}{ejemplo.5.11}{}}
\newlabel{ej:benchmark@cref}{{[ejemplo][11][]5.11}{[1][23][]23}}
\newlabel{eje:5-lote}{{5.1}{24}{}{ejercicio.5.1}{}}
\newlabel{eje:5-lote@cref}{{[ejercicio][1][5]5.1}{[1][24][]24}}
\newlabel{eje:5-vf}{{5.2}{24}{}{ejercicio.5.2}{}}
\newlabel{eje:5-vf@cref}{{[ejercicio][2][5]5.2}{[1][24][]24}}
\newlabel{eje:5-pi}{{5.3}{24}{}{ejercicio.5.3}{}}
\newlabel{eje:5-pi@cref}{{[ejercicio][3][5]5.3}{[1][24][]24}}
\newlabel{eje:5-chebyshev}{{5.4}{24}{}{ejercicio.5.4}{}}
\newlabel{eje:5-chebyshev@cref}{{[ejercicio][4][5]5.4}{[1][24][]24}}
\newlabel{eje:5-comparar}{{5.5}{24}{}{ejercicio.5.5}{}}
\newlabel{eje:5-comparar@cref}{{[ejercicio][5][5]5.5}{[1][24][]24}}
\newlabel{eje:5-n}{{5.6}{24}{}{ejercicio.5.6}{}}
\newlabel{eje:5-n@cref}{{[ejercicio][6][5]5.6}{[1][24][]24}}
\newlabel{ap:soluciones}{{A}{25}{Soluciones de los ejercicios}{appendix.A}{}}
\newlabel{ap:soluciones@cref}{{[appendix][1][2147483647]A}{[1][25][]25}}
\newlabel{sol:eje:1-inversa}{{A.1}{25}{Tema 1}{section*.3}{}}
\newlabel{sol:eje:1-inversa@cref}{{[subappendix][1][2147483647,1]A.1}{[1][25][]25}}
\newlabel{sol:eje:1-disjuntos}{{A.1}{25}{Tema 1}{section*.4}{}}
\newlabel{sol:eje:1-disjuntos@cref}{{[subappendix][1][2147483647,1]A.1}{[1][25][]25}}
\newlabel{sol:eje:1-detector}{{A.1}{26}{Tema 1}{section*.5}{}}
\newlabel{sol:eje:1-detector@cref}{{[subappendix][1][2147483647,1]A.1}{[1][25][]26}}
\newlabel{sol:eje:1-cribado}{{A.1}{26}{Tema 1}{section*.6}{}}
\newlabel{sol:eje:1-cribado@cref}{{[subappendix][1][2147483647,1]A.1}{[1][26][]26}}
\newlabel{sol:eje:1-monedas}{{A.1}{26}{Tema 1}{section*.7}{}}
\newlabel{sol:eje:1-monedas@cref}{{[subappendix][1][2147483647,1]A.1}{[1][26][]26}}
\newlabel{sol:eje:1-complementos}{{A.1}{27}{Tema 1}{section*.8}{}}
\newlabel{sol:eje:1-complementos@cref}{{[subappendix][1][2147483647,1]A.1}{[1][27][]27}}
\newlabel{sol:eje:2-categorica}{{A.2}{27}{Tema 2}{section*.9}{}}
\newlabel{sol:eje:2-categorica@cref}{{[subappendix][2][2147483647,1]A.2}{[1][27][]27}}
\newlabel{sol:eje:2-dados}{{A.2}{28}{Tema 2}{section*.10}{}}
\newlabel{sol:eje:2-dados@cref}{{[subappendix][2][2147483647,1]A.2}{[1][28][]28}}
\newlabel{sol:eje:2-softmax}{{A.2}{28}{Tema 2}{section*.11}{}}
\newlabel{sol:eje:2-softmax@cref}{{[subappendix][2][2147483647,1]A.2}{[1][28][]28}}
\newlabel{sol:eje:2-dados4}{{A.2}{28}{Tema 2}{section*.12}{}}
\newlabel{sol:eje:2-dados4@cref}{{[subappendix][2][2147483647,1]A.2}{[1][28][]28}}
\newlabel{sol:eje:2-binomial}{{A.2}{29}{Tema 2}{section*.13}{}}
\newlabel{sol:eje:2-binomial@cref}{{[subappendix][2][2147483647,1]A.2}{[1][29][]29}}
\newlabel{sol:eje:2-covarianza}{{A.2}{29}{Tema 2}{section*.14}{}}
\newlabel{sol:eje:2-covarianza@cref}{{[subappendix][2][2147483647,1]A.2}{[1][29][]29}}
\newlabel{sol:eje:3-tabla}{{A.3}{30}{Tema 3}{section*.15}{}}
\newlabel{sol:eje:3-tabla@cref}{{[subappendix][3][2147483647,1]A.3}{[1][30][]30}}
\newlabel{sol:eje:3-espera}{{A.3}{30}{Tema 3}{section*.16}{}}
\newlabel{sol:eje:3-espera@cref}{{[subappendix][3][2147483647,1]A.3}{[1][30][]30}}
\newlabel{sol:eje:3-cuadrado}{{A.3}{30}{Tema 3}{section*.17}{}}
\newlabel{sol:eje:3-cuadrado@cref}{{[subappendix][3][2147483647,1]A.3}{[1][30][]30}}
\newlabel{sol:eje:3-init}{{A.3}{31}{Tema 3}{section*.18}{}}
\newlabel{sol:eje:3-init@cref}{{[subappendix][3][2147483647,1]A.3}{[1][31][]31}}
\newlabel{sol:eje:3-exponencial}{{A.3}{31}{Tema 3}{section*.19}{}}
\newlabel{sol:eje:3-exponencial@cref}{{[subappendix][3][2147483647,1]A.3}{[1][31][]31}}
\newlabel{sol:eje:3-he}{{A.3}{31}{Tema 3}{section*.20}{}}
\newlabel{sol:eje:3-he@cref}{{[subappendix][3][2147483647,1]A.3}{[1][31][]31}}
\newlabel{sol:eje:4-ceros}{{A.4}{32}{Tema 4}{section*.21}{}}
\newlabel{sol:eje:4-ceros@cref}{{[subappendix][4][2147483647,1]A.4}{[1][32][]32}}
\newlabel{sol:eje:4-datos}{{A.4}{32}{Tema 4}{section*.22}{}}
\newlabel{sol:eje:4-datos@cref}{{[subappendix][4][2147483647,1]A.4}{[1][32][]32}}
\newlabel{sol:eje:4-moneda}{{A.4}{33}{Tema 4}{section*.23}{}}
\newlabel{sol:eje:4-moneda@cref}{{[subappendix][4][2147483647,1]A.4}{[1][33][]33}}
\newlabel{sol:eje:4-u2}{{A.4}{33}{Tema 4}{section*.24}{}}
\newlabel{sol:eje:4-u2@cref}{{[subappendix][4][2147483647,1]A.4}{[1][33][]33}}
\newlabel{sol:eje:4-u3}{{A.4}{33}{Tema 4}{section*.25}{}}
\newlabel{sol:eje:4-u3@cref}{{[subappendix][4][2147483647,1]A.4}{[1][33][]33}}
\newlabel{sol:eje:4-sesgo}{{A.4}{34}{Tema 4}{section*.26}{}}
\newlabel{sol:eje:4-sesgo@cref}{{[subappendix][4][2147483647,1]A.4}{[1][34][]34}}
\newlabel{sol:eje:4-beta}{{A.4}{34}{Tema 4}{section*.27}{}}
\newlabel{sol:eje:4-beta@cref}{{[subappendix][4][2147483647,1]A.4}{[1][34][]34}}
\newlabel{sol:eje:5-lote}{{A.5}{35}{Tema 5}{section*.28}{}}
\newlabel{sol:eje:5-lote@cref}{{[subappendix][5][2147483647,1]A.5}{[1][35][]35}}
\newlabel{sol:eje:5-vf}{{A.5}{35}{Tema 5}{section*.29}{}}
\newlabel{sol:eje:5-vf@cref}{{[subappendix][5][2147483647,1]A.5}{[1][35][]35}}
\newlabel{sol:eje:5-pi}{{A.5}{35}{Tema 5}{section*.30}{}}
\newlabel{sol:eje:5-pi@cref}{{[subappendix][5][2147483647,1]A.5}{[1][35][]35}}
\newlabel{sol:eje:5-chebyshev}{{A.5}{36}{Tema 5}{section*.31}{}}
\newlabel{sol:eje:5-chebyshev@cref}{{[subappendix][5][2147483647,1]A.5}{[1][36][]36}}
\newlabel{sol:eje:5-comparar}{{A.5}{36}{Tema 5}{section*.32}{}}
\newlabel{sol:eje:5-comparar@cref}{{[subappendix][5][2147483647,1]A.5}{[1][36][]36}}
\newlabel{sol:eje:5-n}{{A.5}{37}{Tema 5}{section*.33}{}}
\newlabel{sol:eje:5-n@cref}{{[subappendix][5][2147483647,1]A.5}{[1][37][]37}}
\makeatother

\begin{document}
\frontmatter
\begin{titlepage}
\centering
\vspace*{2.5cm}
{\large\scshape Anatomía de un LLM}\par
\vspace{0.4cm}
{\normalsize Fundamentos matemáticos de los modelos de lenguaje}\par
\vspace{2.2cm}
\rule{\textwidth}{0.8pt}\par\vspace{0.5cm}
{\small\scshape Fase I \textperiodcentered\ Fundamentos matemáticos}\par\vspace{0.3cm}
{\Huge\bfseries Módulo 5\par}\vspace{0.3cm}
{\LARGE Probabilidad y estadística\par}
\vspace{0.5cm}\rule{\textwidth}{0.4pt}\par
\vspace{1.6cm}
{\large Resumen del módulo}\par
\vfill
{\small Edición de octubre de 2026}\par
\end{titlepage}

\chapter*{Sobre este resumen}
\addcontentsline{toc}{chapter}{Sobre este resumen}
\markboth{Sobre este resumen}{}

Este documento reúne lo indispensable de los apuntes del Módulo 5: el texto expositivo, todas las definiciones y todos los resultados (teoremas, proposiciones, lemas y corolarios), las tablas y las observaciones de contenido. Se han suprimido las demostraciones, los ejemplos, las lecturas, las figuras, los ejercicios y sus soluciones. La numeración de temas, secciones, enunciados y ecuaciones es la misma que en los apuntes, de modo que cualquier referencia puede buscarse allí; las que remiten a elementos suprimidos (un ejemplo, una figura, una ecuación de una demostración) conservan su número en los apuntes.

\section*{Notación}
\begin{center}
\begin{tabular}{@{}ll@{}}
\toprule
Símbolo & Significado \\
\midrule
$(\Omega,\mathcal{F},\PP)$ & espacio de probabilidad \\
$X$, $x$ & variable aleatoria y valor que toma \\
$p_X(x)$, $f_X(x)$, $F_X(x)$ & función de masa, densidad y función de distribución \\
$\E[X]$, $\Var(X)$, $\Cov(X,Y)$ & esperanza, varianza y covarianza \\
$\Ber(p)$, $\Bin(n,p)$, $\Cat(p_1,\dots,p_K)$ & distribuciones de Bernoulli, binomial y categórica \\
$\Unif(a,b)$, $\Expo(\lambda)$, $\N(\mu,\sigma^{2})$ & distribuciones uniforme, exponencial y normal \\
$\Phi$ & función de distribución de la normal estándar \\
$\BetaD(a,b)$, $\Dir(\alpha_1,\dots,\alpha_K)$ & distribuciones beta y de Dirichlet \\
$X_1,\dots,X_n$ \iid & muestra independiente e idénticamente distribuida \\
$\Xbar_n$ & media muestral $\frac1n\sum_{i=1}^{n}X_i$ \\
$\ell(\theta)$, $\hat\theta$ & log-verosimilitud y estimador \\
$\mathbf{1}_A$ & función indicadora del suceso $A$ \\
\bottomrule
\end{tabular}
\end{center}


\tableofcontents

\mainmatter
% =====================================================================
\chapter{Probabilidad, condicionamiento y teorema de Bayes}
\label{tema:1}
% =====================================================================

Este tema fija el lenguaje del módulo: espacios de probabilidad, probabilidad condicionada, independencia y teorema de Bayes. La presentación sigue el capítulo~1 de~\cite{bt} y las lecciones~1 a~3 de~\cite{mit6041}; la terminología bayesiana (a priori, verosimilitud, a posteriori) y la forma de razón de probabilidades siguen las clases~3 y~12 de~\cite{mit1805} y el capítulo~2 de~\cite{bh}. Deisenroth \emph{et al.}~\cite[§6.1 y §6.3]{mml} ofrecen el mismo material con orientación al aprendizaje automático.

\section{Modelos probabilísticos}

\begin{definicion}[Espacio de probabilidad]\label{def:espacio}
Un \emph{espacio de probabilidad} es una terna $(\Omega,\mathcal{F},\PP)$ formada por un \emph{espacio muestral} $\Omega$ (el conjunto de resultados posibles de un experimento), una familia $\mathcal{F}$ de subconjuntos de $\Omega$, llamados \emph{sucesos}, cerrada por complementos y uniones numerables y que contiene a $\Omega$, y una función $\PP\colon\mathcal{F}\to\R$ que cumple los \emph{axiomas de Kolmogorov}:
\begin{enumerate}
  \item \emph{no negatividad}: $\PP(A)\ge 0$ para todo $A\in\mathcal{F}$;
  \item \emph{normalización}: $\PP(\Omega)=1$;
  \item \emph{aditividad numerable}: si $A_1,A_2,\dots$ son disjuntos dos a dos, $\PP\bigl(\bigcup_{i}A_i\bigr)=\sum_{i}\PP(A_i)$.
\end{enumerate}
\end{definicion}

Cuando $\Omega$ es finito o numerable se toma $\mathcal{F}$ igual a todos los subconjuntos de $\Omega$, y basta dar la probabilidad de cada resultado: $\PP(A)=\sum_{\omega\in A}\PP(\{\omega\})$. Si además $\Omega$ es finito y todos los resultados son equiprobables, $\PP(A)=\abs{A}/\abs{\Omega}$ (ley de Laplace)~\cite[§1.2]{bt}.

\begin{proposicion}[Consecuencias de los axiomas]\label{prop:axiomas}
Para sucesos $A,B\in\mathcal{F}$:
\begin{enumerate}
  \item $\PP(\varnothing)=0$ y $\PP(A^{c})=1-\PP(A)$;
  \item si $A\subseteq B$, entonces $\PP(A)\le\PP(B)$;
  \item $\PP(A\cup B)=\PP(A)+\PP(B)-\PP(A\cap B)$;
  \item $\PP(A\cup B)\le\PP(A)+\PP(B)$ (cota de la unión).
\end{enumerate}
\end{proposicion}

\section{Probabilidad condicionada}

\begin{definicion}[Probabilidad condicionada]
Sean $A,B\in\mathcal{F}$ con $\PP(A)>0$. La \emph{probabilidad de $B$ condicionada a $A$} es
\[
\PP(B\mid A)=\frac{\PP(A\cap B)}{\PP(A)}.
\]
\end{definicion}

Condicionar a $A$ equivale a tomar $A$ como nuevo espacio muestral y renormalizar. Para $A$ fijo, $B\mapsto\PP(B\mid A)$ cumple los tres axiomas de Kolmogorov~\cite[§1.3]{bt}; en consecuencia, todas las propiedades de la \cref{prop:axiomas} valen también para probabilidades condicionadas. Despejando se obtiene la \emph{regla del producto}, $\PP(A\cap B)=\PP(A)\,\PP(B\mid A)$, cuya iteración es la regla de la cadena del módulo~1.

\begin{observacion}
Un modelo de lenguaje es una familia de probabilidades condicionadas $\PP(x_t\mid x_{1:t-1})$: el suceso condicionante es el contexto observado y la distribución condicionada es la del token siguiente. Todo lo que se demuestre en este tema para probabilidades condicionadas se aplica, por tanto, a las salidas de un LLM.
\end{observacion}

\section{Probabilidad total y regla de Bayes}

\begin{definicion}[Partición]
Una familia finita o numerable de sucesos $H_1,H_2,\dots$ es una \emph{partición} de $\Omega$ si son disjuntos dos a dos y su unión es $\Omega$.
\end{definicion}

\begin{teorema}[Probabilidad total]\label{teo:total}
Sea $H_1,H_2,\dots$ una partición de $\Omega$ con $\PP(H_i)>0$ para todo $i$. Para todo suceso $D$,
\[
\PP(D)=\sum_{i}\PP(D\mid H_i)\,\PP(H_i).
\]
\end{teorema}

\begin{teorema}[Regla de Bayes]\label{teo:bayes}
Con las hipótesis del \cref{teo:total} y $\PP(D)>0$, para cada $j$,
\[
\PP(H_j\mid D)=\frac{\PP(D\mid H_j)\,\PP(H_j)}{\sum_{i}\PP(D\mid H_i)\,\PP(H_i)}.
\]
\end{teorema}

En el lenguaje de la inferencia, los $H_i$ son \emph{hipótesis} y $D$ son los \emph{datos}. $\PP(H_j)$ es la probabilidad \emph{a priori}, $\PP(D\mid H_j)$ la \emph{verosimilitud}, $\PP(H_j\mid D)$ la probabilidad \emph{a posteriori} y el denominador, $\PP(D)$, la \emph{evidencia}, que solo normaliza. Abreviadamente, la probabilidad a posteriori es proporcional al producto de verosimilitud y probabilidad a priori~\cite[clase~11]{mit1805}.

\addtocounter{definicion}{1}%

\addtocounter{figure}{1}%

Cuando solo hay dos hipótesis, $H$ y $H^{c}$, es cómodo trabajar con la \emph{razón de probabilidades} (en inglés, \emph{odds}) $O(H)=\PP(H)/\PP(H^{c})$.

\begin{proposicion}[Bayes en forma de razón]\label{prop:odds}
Si $\PP(D\mid H^{c})>0$, entonces
\[
O(H\mid D)=\frac{\PP(D\mid H)}{\PP(D\mid H^{c})}\cdot O(H),
\]
donde $O(H\mid D)=\PP(H\mid D)/\PP(H^{c}\mid D)$. El cociente $\PP(D\mid H)/\PP(D\mid H^{c})$ se llama \emph{factor de Bayes} o razón de verosimilitudes.
\end{proposicion}

La forma de razón muestra con claridad la llamada \emph{falacia de la probabilidad a priori}: un factor de Bayes grande no basta para que la hipótesis sea probable si $O(H)$ es muy pequeña (\cref{eje:1-detector,eje:1-cribado}).

\section{Independencia}

\begin{definicion}[Independencia]\label{def:indep}
Dos sucesos $A$ y $B$ son \emph{independientes} si $\PP(A\cap B)=\PP(A)\,\PP(B)$. Si $\PP(A)>0$, la condición equivale a $\PP(B\mid A)=\PP(B)$: saber que ha ocurrido $A$ no cambia la probabilidad de $B$. Los sucesos $A_1,\dots,A_n$ son \emph{independientes} si $\PP\bigl(\bigcap_{i\in S}A_i\bigr)=\prod_{i\in S}\PP(A_i)$ para todo subconjunto $S\subseteq\{1,\dots,n\}$.
\end{definicion}

\begin{definicion}[Independencia condicionada]
Sea $C$ con $\PP(C)>0$. Los sucesos $A$ y $B$ son \emph{independientes condicionalmente a $C$} si $\PP(A\cap B\mid C)=\PP(A\mid C)\,\PP(B\mid C)$.
\end{definicion}

\begin{observacion}
La independencia condicionada no implica la independencia, ni al revés. Dos lanzamientos de una moneda elegida al azar entre una equilibrada y una trucada son independientes condicionalmente a la moneda, pero no lo son incondicionalmente: una cara en el primer lanzamiento hace más probable la moneda que favorece la cara y, con ella, una cara en el segundo (\cref{eje:1-monedas}). Del mismo modo, independencia no es lo mismo que incompatibilidad: dos sucesos disjuntos con probabilidad positiva nunca son independientes (\cref{eje:1-disjuntos}).
\end{observacion}

\addtocounter{definicion}{1}%


% =====================================================================
\chapter{Variables aleatorias discretas}
\label{tema:2}
% =====================================================================

Una variable aleatoria asigna un número a cada resultado de un experimento. Este tema estudia las variables discretas, entre ellas la distribución categórica que produce un LLM en cada paso, y sus dos resúmenes principales, la esperanza y la varianza. Se sigue el capítulo~2 y la §4.2 de~\cite{bt}, las lecciones~5 a~7 de~\cite{mit6041} y las clases~4, 5a y~7 de~\cite{mit1805}; la orientación al aprendizaje automático procede de~\cite[§6.2 y §6.4]{mml}.

\section{Variables aleatorias y funciones de masa}

\begin{definicion}[Variable aleatoria discreta]\label{def:va}
Una \emph{variable aleatoria} es una función $X\colon\Omega\to\R$ tal que $\{\omega: X(\omega)\le x\}\in\mathcal{F}$ para todo $x\in\R$. Es \emph{discreta} si toma un conjunto finito o numerable de valores. Su \emph{función de masa} es
\[
p_X(x)=\PP(X=x),
\]
que cumple $p_X(x)\ge 0$ y $\sum_x p_X(x)=1$. Si $X$ y $Y$ tienen la misma función de masa se dice que tienen la misma \emph{distribución}.
\end{definicion}

\begin{definicion}[Distribuciones discretas básicas]\label{def:discretas}
\begin{enumerate}
  \item $X\sim\Ber(p)$, con $p\in[0,1]$, si $X\in\{0,1\}$ y $p_X(x)=p^{x}(1-p)^{1-x}$. Modela un experimento con dos resultados: acierto o fallo.
  \item $X\sim\Bin(n,p)$ si $X$ es el número de éxitos en $n$ ensayos de Bernoulli independientes con probabilidad $p$:
  \[
  p_X(k)=\binom{n}{k}p^{k}(1-p)^{n-k},\qquad k=0,1,\dots,n.
  \]
  \item $X\sim\Cat(p_1,\dots,p_K)$, con $p_k\ge 0$ y $\sum_k p_k=1$, si $X\in\{1,\dots,K\}$ y $\PP(X=k)=p_k$. Equivalentemente, $p_X(k)=\prod_{j=1}^{K}p_j^{\mathbf{1}[j=k]}$.
\end{enumerate}
\end{definicion}

La categórica es la generalización de la Bernoulli a $K$ resultados, y es exactamente lo que produce un LLM en cada posición: la softmax convierte los $K$ logits en el vector $(p_1,\dots,p_K)$ de una categórica sobre el vocabulario. Los valores $1,\dots,K$ son etiquetas sin orden ni distancia: el token~517 no está «entre» el~516 y el~518.

\addtocounter{definicion}{1}%

\section{Esperanza}

\begin{definicion}[Esperanza]\label{def:esperanza}
La \emph{esperanza} o \emph{media} de una variable discreta $X$ es
\[
\E[X]=\sum_x x\,p_X(x),
\]
siempre que la serie converja absolutamente.
\end{definicion}

La esperanza es el centro de gravedad de la función de masa. Para una categórica con etiquetas arbitrarias no tiene sentido «el token medio», pero sí la esperanza de cualquier función numérica de los tokens, por ejemplo de la pérdida.

\begin{teorema}[Esperanza de una función]\label{teo:lotus}
Si $g\colon\R\to\R$ y $X$ es discreta, entonces $\E[g(X)]=\sum_x g(x)\,p_X(x)$.
\end{teorema}

\begin{proposicion}[Linealidad]\label{prop:linealidad}
Para variables $X$, $Y$ con esperanza finita y constantes $a,b,c$,
\[
\E[aX+bY+c]=a\,\E[X]+b\,\E[Y]+c.
\]
La igualdad no requiere independencia.
\end{proposicion}

\addtocounter{definicion}{1}%

\section{Varianza}

\begin{definicion}[Varianza]\label{def:varianza}
La \emph{varianza} de $X$ es $\Var(X)=\E\bigl[(X-\E[X])^{2}\bigr]$, y su \emph{desviación típica} es $\sigma_X=\sqrt{\Var(X)}$, que tiene las mismas unidades que $X$.
\end{definicion}

\begin{proposicion}\label{prop:varianza}
\begin{enumerate}
  \item $\Var(X)=\E[X^{2}]-\E[X]^{2}$.
  \item $\Var(aX+b)=a^{2}\Var(X)$ para constantes $a,b$.
  \item $\Var(X)\ge 0$, con igualdad si y solo si $X$ es constante con probabilidad~1.
\end{enumerate}
\end{proposicion}

\addtocounter{definicion}{1}%

Para $X\sim\Ber(p)$, $X^{2}=X$, luego $\E[X]=\E[X^{2}]=p$ y $\Var(X)=p-p^{2}=p(1-p)$, máxima en $p=1/2$.

\section{Varias variables aleatorias}

\begin{definicion}[Distribución conjunta e independencia]\label{def:conjunta}
La \emph{función de masa conjunta} de $X$ e $Y$ es $p_{X,Y}(x,y)=\PP(X=x,\,Y=y)$; las \emph{marginales} se obtienen sumando, $p_X(x)=\sum_y p_{X,Y}(x,y)$. Las variables $X$ e $Y$ son \emph{independientes} si $p_{X,Y}(x,y)=p_X(x)\,p_Y(y)$ para todo $(x,y)$, y $X_1,\dots,X_n$ son independientes si la conjunta factoriza en el producto de las $n$ marginales.
\end{definicion}

\begin{definicion}[Covarianza]
La \emph{covarianza} de $X$ e $Y$ es $\Cov(X,Y)=\E\bigl[(X-\E[X])(Y-\E[Y])\bigr]=\E[XY]-\E[X]\,\E[Y]$. Si $\Cov(X,Y)=0$, las variables se llaman \emph{incorreladas}.
\end{definicion}

\begin{proposicion}\label{prop:suma}
\begin{enumerate}
  \item Si $X$ e $Y$ son independientes, $\E[XY]=\E[X]\,\E[Y]$ y, por tanto, $\Cov(X,Y)=0$.
  \item $\Var(X+Y)=\Var(X)+\Var(Y)+2\Cov(X,Y)$. En particular, la varianza de una suma de variables independientes es la suma de las varianzas.
\end{enumerate}
\end{proposicion}

El recíproco de (i) es falso: hay variables incorreladas que no son independientes (\cref{eje:2-covarianza}).

\begin{corolario}[Varianza de la media muestral]\label{cor:media}
Sean $X_1,\dots,X_n$ independientes con media $\mu$ y varianza $\sigma^{2}$, y sea $\Xbar_n=\frac1n\sum_{i=1}^{n}X_i$. Entonces $\E[\Xbar_n]=\mu$ y $\Var(\Xbar_n)=\sigma^{2}/n$.
\end{corolario}

El \cref{cor:media} es el resultado central del módulo: promediar $n$ observaciones independientes divide la varianza por $n$ y la desviación típica por $\sqrt n$. El tema~5 lo convierte en garantías sobre el error de cualquier promedio. Como aplicación inmediata, si $X\sim\Bin(n,p)$ se escribe $X=\sum_{i=1}^{n}B_i$ con $B_i\sim\Ber(p)$ independientes, y se obtiene $\E[X]=np$ y $\Var(X)=np(1-p)$.

\begin{table}[h]
\centering\small
\caption{Distribuciones de los temas~2 y~3 y uso en un LLM.}\label{tab:distribuciones}
\begin{tabular}{@{}lllll@{}}
\toprule
Distribución & Valores & $\E[X]$ & $\Var(X)$ & Uso en un LLM \\
\midrule
$\Ber(p)$ & $\{0,1\}$ & $p$ & $p(1-p)$ & acierto en una pregunta de evaluación \\
$\Bin(n,p)$ & $\{0,\dots,n\}$ & $np$ & $np(1-p)$ & aciertos en un banco de pruebas \\
$\Cat(p_1,\dots,p_K)$ & $\{1,\dots,K\}$ & --- & --- & token siguiente \\
$\Unif(a,b)$ & $[a,b]$ & $(a+b)/2$ & $(b-a)^{2}/12$ & número aleatorio del muestreo \\
$\Expo(\lambda)$ & $[0,\infty)$ & $1/\lambda$ & $1/\lambda^{2}$ & tiempos de espera \\
$\N(\mu,\sigma^{2})$ & $\R$ & $\mu$ & $\sigma^{2}$ & pesos iniciales, ruido \\
\bottomrule
\end{tabular}
\end{table}

\section{Muestreo de una distribución categórica}

Para generar un token, el LLM necesita extraer una muestra de $\Cat(p_1,\dots,p_K)$ a partir de un número aleatorio uniforme. Sea $F_k=p_1+\dots+p_k$ la probabilidad acumulada, con $F_0=0$ y $F_K=1$.

\begin{proposicion}[Transformada inversa discreta]\label{prop:inversa-discreta}
Sea $U\sim\Unif(0,1)$, es decir, $\PP(a<U\le b)=b-a$ para $0\le a\le b\le 1$ (\cref{tema:3}). La variable
\[
X=\min\{k: U\le F_k\}
\]
tiene distribución $\Cat(p_1,\dots,p_K)$.
\end{proposicion}

\addtocounter{figure}{1}%

\addtocounter{definicion}{1}%

\addtocounter{definicion}{1}%


% =====================================================================
\chapter{Variables aleatorias continuas y distribución normal}
\label{tema:3}
% =====================================================================

Las magnitudes que varían en un continuo, como los pesos de una red o el número aleatorio que se usa para muestrear, se describen con densidades. Este tema introduce las variables continuas, estudia la uniforme, la exponencial y la normal, demuestra el método de la transformada inversa y aplica las propiedades de la varianza a la inicialización de redes neuronales. Se sigue el capítulo~3 y la §4.1 de~\cite{bt}, la unidad~II de~\cite{mit6041}, las clases~5b a~6a de~\cite{mit1805} y el capítulo~5 de~\cite{bh}; la normal en aprendizaje automático, en~\cite[§6.5]{mml} y~\cite[§1.2.4]{bishop}, y la inicialización, en los artículos originales~\cite{glorot,he}.

\section{Densidad y función de distribución}

\begin{definicion}[Variable continua]\label{def:continua}
Una variable aleatoria $X$ es \emph{continua} si existe una función $f_X\ge 0$, su \emph{densidad}, tal que
\[
\PP(a\le X\le b)=\int_a^b f_X(x)\,dx\qquad\text{para todo } a\le b,
\]
y en particular $\int_{-\infty}^{\infty}f_X(x)\,dx=1$. La \emph{función de distribución} de cualquier variable es $F_X(x)=\PP(X\le x)$.
\end{definicion}

\begin{proposicion}\label{prop:densidad}
Si $X$ es continua con densidad $f_X$:
\begin{enumerate}
  \item $\PP(X=x)=0$ para todo $x$; por tanto, da igual usar desigualdades estrictas o no estrictas;
  \item $F_X(x)=\int_{-\infty}^{x}f_X(t)\,dt$ y, en cada punto de continuidad de $f_X$, $F_X'(x)=f_X(x)$;
  \item $F_X$ es no decreciente, con $\lim_{x\to-\infty}F_X(x)=0$ y $\lim_{x\to\infty}F_X(x)=1$.
\end{enumerate}
\end{proposicion}

La densidad no es una probabilidad: puede ser mayor que~1. Para $h$ pequeño, $\PP(x\le X\le x+h)\approx f_X(x)\,h$, de modo que $f_X(x)$ es probabilidad por unidad de longitud. La esperanza y la varianza se definen como en el caso discreto cambiando sumas por integrales: $\E[g(X)]=\int g(x)f_X(x)\,dx$, y todas las propiedades del \cref{tema:2} se mantienen~\cite[§3.1]{bt}.

\section{Uniforme y exponencial}

\begin{definicion}[Uniforme]
$X\sim\Unif(a,b)$, con $a<b$, si $f_X(x)=1/(b-a)$ en $[a,b]$ y $0$ fuera.
\end{definicion}

\begin{proposicion}\label{prop:uniforme}
Si $X\sim\Unif(a,b)$, entonces $\E[X]=(a+b)/2$ y $\Var(X)=(b-a)^{2}/12$.
\end{proposicion}

\begin{definicion}[Exponencial]
$X\sim\Expo(\lambda)$, con $\lambda>0$, si $f_X(x)=\lambda e^{-\lambda x}$ para $x\ge 0$ y $0$ para $x<0$. Su función de distribución es $F_X(x)=1-e^{-\lambda x}$ para $x\ge 0$.
\end{definicion}

Integrando por partes se obtiene $\E[X]=1/\lambda$ y $\E[X^{2}]=2/\lambda^{2}$, luego $\Var(X)=1/\lambda^{2}$~\cite[§3.1]{bt}. La exponencial modela tiempos de espera sin desgaste, como el de un taxi o el de una desintegración radiactiva~\cite[clase~5c]{mit1805}.

\begin{proposicion}[Falta de memoria]\label{prop:memoria}
Si $X\sim\Expo(\lambda)$ y $s,t\ge 0$, entonces $\PP(X>s+t\mid X>s)=\PP(X>t)$.
\end{proposicion}

\section{La distribución normal}

\begin{definicion}[Normal]\label{def:normal}
$X\sim\N(\mu,\sigma^{2})$, con $\mu\in\R$ y $\sigma>0$, si
\[
f_X(x)=\frac{1}{\sqrt{2\pi\sigma^{2}}}\exp\Bigl(-\frac{(x-\mu)^{2}}{2\sigma^{2}}\Bigr).
\]
La \emph{normal estándar} es $\N(0,1)$; su función de distribución se denota $\Phi$, y el valor máximo de su densidad es $1/\sqrt{2\pi}\approx 0{,}399$.
\end{definicion}

Que la densidad integre~1 se deduce de la integral gaussiana $\int_{-\infty}^{\infty}e^{-x^{2}/2}\,dx=\sqrt{2\pi}$, que se calcula pasando a coordenadas polares en el plano~\cite[§3.3]{bt}. Por simetría, $\E[X]=\mu$, y una integración por partes da $\Var(X)=\sigma^{2}$.

\begin{proposicion}[Transformaciones lineales]\label{prop:afin}
Si $X\sim\N(\mu,\sigma^{2})$ y $a\ne 0$, $b$ son constantes, entonces $aX+b\sim\N(a\mu+b,\,a^{2}\sigma^{2})$. En particular, $Z=(X-\mu)/\sigma\sim\N(0,1)$ y $\PP(X\le x)=\Phi\bigl((x-\mu)/\sigma\bigr)$.
\end{proposicion}

Gracias a la \cref{prop:afin}, cualquier probabilidad normal se reduce a valores de $\Phi$, que se tabulan o se calculan numéricamente. El \cref{tab:normal} recoge los más usados y la \cref{fig:normal} los representa.

\begin{table}[h]
\centering\small
\caption{Probabilidad de que una normal caiga a menos de $k$ desviaciones típicas de su media.}\label{tab:normal}
\begin{tabular}{@{}lcccc@{}}
\toprule
$k$ & $1$ & $1{,}96$ & $2$ & $3$ \\
\midrule
$\PP(\abs{X-\mu}\le k\sigma)=2\Phi(k)-1$ & $0{,}683$ & $0{,}950$ & $0{,}954$ & $0{,}997$ \\
\bottomrule
\end{tabular}
\end{table}

\addtocounter{figure}{1}%

\addtocounter{definicion}{1}%

La normal aparece cuando una magnitud es la suma de muchos efectos pequeños e independientes; el teorema central del límite (\cref{tema:5}) da la justificación precisa.

\section{Muestreo por transformada inversa}

Un ordenador genera números (pseudo)aleatorios uniformes en $(0,1)$. El resultado siguiente permite convertirlos en muestras de cualquier distribución continua cuya función de distribución sea invertible; es la versión continua de la \cref{prop:inversa-discreta}.

\begin{teorema}[Universalidad de la uniforme]\label{teo:inversa}
Sea $F$ una función de distribución continua y estrictamente creciente en el intervalo donde $0<F<1$, y sea $U\sim\Unif(0,1)$. Entonces $X=F^{-1}(U)$ tiene función de distribución $F$.
\end{teorema}

\begin{observacion}
El \cref{teo:inversa} es la base de la generación de muestras en todas las bibliotecas numéricas. Para la normal, $\Phi^{-1}$ no tiene expresión cerrada y se usan otros métodos (por ejemplo, el de Box--Muller), pero el principio es el mismo: transformar uniformes~\cite[§6.7]{mml}. El \cref{eje:3-exponencial} lo aplica a la exponencial.
\end{observacion}

\section{Aplicación: inicialización de redes neuronales}

Una neurona lineal calcula $y=\sum_{j=1}^{n}w_jx_j$, donde $n$ es el número de entradas (el \emph{fan-in}). Antes de entrenar, los pesos $w_j$ se eligen al azar. ¿Con qué varianza?

\begin{proposicion}[Varianza de una neurona]\label{prop:neurona}
Supóngase que $w_1,\dots,w_n,x_1,\dots,x_n$ son independientes, que los $w_j$ tienen media~0 y varianza $\Var(w)$ y que los $x_j$ tienen media~0 y varianza $\Var(x)$. Entonces
\[
\E[y]=0,\qquad \Var(y)=n\,\Var(w)\,\Var(x).
\]
\end{proposicion}

Si se toma $\Var(w)=1$ y $n=512$, cada capa multiplica la varianza por~512 y la desviación típica por $\sqrt{512}\approx 22{,}6$; tras diez capas la desviación ha crecido en un factor $512^{5}\approx 3{,}5\cdot 10^{13}$ y la red se desborda. Si $\Var(w)$ es demasiado pequeña, la señal se anula. La \cref{prop:neurona} da la elección que conserva la varianza:
\[
\Var(w)=\frac1n,\qquad\text{por ejemplo } n=512\ \Rightarrow\ \sigma_w=1/\sqrt{512}\approx 0{,}0442.
\]
Glorot y Bengio~\cite{glorot} propusieron este análisis y, para equilibrar también la propagación hacia atrás de los gradientes, la elección de compromiso $\Var(w)=2/(n_{\mathrm{ent}}+n_{\mathrm{sal}})$. Con la activación ReLU, $\max(0,x)$, la mitad de las entradas se anula y el segundo momento se reduce a la mitad (\cref{eje:3-he}); He \emph{et al.}~\cite{he} compensan con $\Var(w)=2/n$, que para $n=512$ da $\sigma_w=\sqrt{2/512}=0{,}0625$. El módulo~9 retoma estas reglas.


% =====================================================================
\chapter{Estimación de parámetros}
\label{tema:4}
% =====================================================================

Entrenar un modelo es elegir sus parámetros a partir de datos. Este tema presenta los dos enfoques clásicos de la estimación: el frecuentista, cuyo método principal es la máxima verosimilitud, y el bayesiano, que trata los parámetros como variables aleatorias con una distribución a priori. Se sigue la §9.1 y las §8.1–8.2 de~\cite{bt}, las lecciones~21 a~23 de~\cite{mit6041} y las clases~10 a~15 de~\cite{mit1805}; la conexión con el aprendizaje automático procede de~\cite[§6.6 y §8.3]{mml} y de~\cite[§2.1–2.3]{bishop}.

\section{Estimadores, sesgo y error cuadrático medio}

\begin{definicion}[Modelo estadístico y estimador]
Un \emph{modelo estadístico} es una familia de distribuciones $\{p(\cdot\,;\theta):\theta\in\Theta\}$ indexada por un parámetro $\theta$. Dada una muestra $X_1,\dots,X_n$ \iid\ con distribución $p(\cdot\,;\theta)$ para un $\theta$ desconocido, un \emph{estimador} es una función $\hat\theta_n=g(X_1,\dots,X_n)$ de la muestra. Es una variable aleatoria.
\end{definicion}

\begin{definicion}[Sesgo y error cuadrático medio]
El \emph{sesgo} de $\hat\theta_n$ es $b(\hat\theta_n)=\E[\hat\theta_n]-\theta$; el estimador es \emph{insesgado} si $b=0$ para todo $\theta$. Su \emph{error cuadrático medio} es $\mathrm{ECM}(\hat\theta_n)=\E\bigl[(\hat\theta_n-\theta)^{2}\bigr]$~\cite[§9.1]{bt}.
\end{definicion}

\begin{proposicion}[Descomposición sesgo--varianza]\label{prop:ecm}
$\mathrm{ECM}(\hat\theta_n)=b(\hat\theta_n)^{2}+\Var(\hat\theta_n)$.
\end{proposicion}

Por el \cref{cor:media}, la media muestral $\Xbar_n$ es un estimador insesgado de la media, con $\mathrm{ECM}=\sigma^{2}/n$.

\section{Máxima verosimilitud}

\begin{definicion}[Estimador de máxima verosimilitud]\label{def:emv}
Dada una realización $x_1,\dots,x_n$ de una muestra \iid, la \emph{verosimilitud} es $\mathcal{L}(\theta)=\prod_{i=1}^{n}p(x_i;\theta)$ y la \emph{log-verosimilitud} es $\ell(\theta)=\sum_{i=1}^{n}\ln p(x_i;\theta)$. El \emph{estimador de máxima verosimilitud} (EMV) es
\[
\hat\theta_{\mathrm{EMV}}=\argmax_{\theta\in\Theta}\ell(\theta)=\argmin_{\theta\in\Theta}\Bigl(-\frac1n\sum_{i=1}^{n}\ln p(x_i;\theta)\Bigr).
\]
\end{definicion}

Las dos formas son equivalentes porque $\ln$ es creciente y dividir por $-n$ invierte el sentido de la optimización. La segunda es la log-verosimilitud negativa media, que el módulo~1 identificó con la entropía cruzada entre la distribución empírica de los datos y el modelo. En consecuencia, \emph{preentrenar un LLM minimizando la entropía cruzada es calcular un estimador de máxima verosimilitud}, con $\theta$ de miles de millones de componentes~\cite[§8.3]{mml}.

\addtocounter{definicion}{1}%

Para la categórica con conteos $N_1,\dots,N_K$, el EMV es la frecuencia relativa $\hat p_j=N_j/\sum_k N_k$; es el teorema del EMV del modelo de bigramas del módulo~1, demostrado allí con multiplicadores de Lagrange. Contar y normalizar es, por tanto, máxima verosimilitud.

\section{El EMV de la normal}

\begin{teorema}[EMV de la normal]\label{teo:emv-normal}
Sea $x_1,\dots,x_n$ una muestra de $\N(\mu,\sigma^{2})$ con $n\ge 2$ y no todos los $x_i$ iguales. El EMV de $(\mu,\sigma^{2})$ es
\[
\hat\mu=\bar x=\frac1n\sum_{i=1}^{n}x_i,\qquad \hat\sigma^{2}=\frac1n\sum_{i=1}^{n}(x_i-\bar x)^{2}.
\]
\end{teorema}

\addtocounter{definicion}{1}%

\begin{proposicion}[Sesgo de la varianza muestral]\label{prop:sesgo}
Si $X_1,\dots,X_n$ son \iid\ con varianza $\sigma^{2}$, entonces $\E[\hat\sigma^{2}]=\frac{n-1}{n}\sigma^{2}$. En consecuencia, $s^{2}=\frac{1}{n-1}\sum_i(X_i-\Xbar_n)^{2}$ es un estimador insesgado de $\sigma^{2}$.
\end{proposicion}

En el \cref{ej:emv-normal}, la versión insesgada es $s^{2}=32/7\approx 4{,}571$. El sesgo se debe a que $\bar x$ se ajusta a los propios datos y las desviaciones respecto de él son, en media, menores que respecto de $\mu$. El EMV no es siempre insesgado, pero su sesgo tiende a~0 cuando $n\to\infty$.

\begin{proposicion}[Mínimos cuadrados como máxima verosimilitud]\label{prop:mco}
Supóngase $y_i=f_\theta(x_i)+\varepsilon_i$ con $\varepsilon_i\sim\N(0,\sigma^{2})$ independientes y $\sigma$ fija. Maximizar la verosimilitud de $(y_1,\dots,y_n)$ en $\theta$ equivale a minimizar el error cuadrático $\sum_i\bigl(y_i-f_\theta(x_i)\bigr)^{2}$.
\end{proposicion}

El módulo~8 usa la \cref{prop:mco} para justificar la regresión lineal; la entropía cruzada de la clasificación es su análogo para la categórica.

\section{Inferencia bayesiana y distribuciones conjugadas}

El EMV de la categórica asigna probabilidad nula a toda categoría no observada, y la pérdida de un texto que la contenga es infinita (módulo~1). El enfoque bayesiano lo remedia tratando el parámetro como una variable aleatoria con una distribución \emph{a priori} $\pi(\theta)$ y actualizándola con la regla de Bayes~\cite[§8.1]{bt}:
\[
\pi(\theta\mid x_{1:n})=\frac{\mathcal{L}(\theta)\,\pi(\theta)}{\int\mathcal{L}(\theta')\,\pi(\theta')\,d\theta'}\ \propto\ \mathcal{L}(\theta)\,\pi(\theta).
\]

\begin{definicion}[Estimadores bayesianos]
La \emph{media a posteriori} es $\E[\theta\mid x_{1:n}]$. El estimador \emph{máximo a posteriori} (MAP) es $\hat\theta_{\mathrm{MAP}}=\argmax_\theta\pi(\theta\mid x_{1:n})=\argmax_\theta\bigl(\ell(\theta)+\ln\pi(\theta)\bigr)$~\cite[§8.2]{bt}.
\end{definicion}

\begin{definicion}[Distribución beta]
$p\sim\BetaD(a,b)$, con $a,b>0$, si su densidad en $(0,1)$ es proporcional a $p^{a-1}(1-p)^{b-1}$. Su media es $a/(a+b)$ y, si $a,b>1$, su moda es $(a-1)/(a+b-2)$~\cite[§2.1.1]{bishop}.
\end{definicion}

\begin{teorema}[Conjugación beta--Bernoulli]\label{teo:beta}
Si $p\sim\BetaD(a,b)$ a priori y se observan $k$ éxitos en $n$ ensayos de Bernoulli independientes dado $p$, la distribución a posteriori es $\BetaD(a+k,\,b+n-k)$. En particular,
\[
\E[p\mid\text{datos}]=\frac{k+a}{n+a+b}.
\]
\end{teorema}

Se dice que la beta es \emph{conjugada} de la Bernoulli: la distribución a posteriori pertenece a la misma familia que la a priori, y los datos se limitan a sumar conteos a los parámetros, que actúan como \emph{pseudoconteos}~\cite[clase~15]{mit1805}. El mismo argumento vale para $K$ categorías con la distribución de Dirichlet, de densidad proporcional a $\prod_{j=1}^{K}p_j^{\alpha_j-1}$ sobre el símplex: con conteos $N_1,\dots,N_K$, la distribución a posteriori es $\Dir(\alpha_1+N_1,\dots,\alpha_K+N_K)$~\cite[§2.2.1]{bishop}, \cite[§6.6]{mml}.

\begin{corolario}[Suavizado de Laplace como estimación bayesiana]\label{cor:laplace}
Con distribución a priori $\Dir(\alpha,\dots,\alpha)$ y conteos $N_1,\dots,N_K$ de total $N$,
\[
\E[p_j\mid\text{datos}]=\frac{N_j+\alpha}{N+K\alpha},\qquad
\hat p_{j,\mathrm{MAP}}=\frac{N_j+\alpha-1}{N+K(\alpha-1)}\quad(\alpha>1).
\]
El suavizado de Laplace del módulo~1 es la media a posteriori con $\alpha=1$ (a priori uniforme sobre el símplex) y también el MAP con $\alpha=2$.
\end{corolario}

\begin{table}[h]
\centering\small
\caption{Estimadores de una categórica con conteos $N=(3,1,0)$, $K=3$.}\label{tab:estimadores}
\begin{tabular}{@{}lll@{}}
\toprule
Estimador & Fórmula para $p_j$ & Resultado \\
\midrule
EMV & $N_j/N$ & $(0{,}75;\ 0{,}25;\ 0)$ \\
Media a posteriori, $\alpha=1$ & $(N_j+1)/(N+K)$ & $(4/7;\ 2/7;\ 1/7)\approx(0{,}571;\ 0{,}286;\ 0{,}143)$ \\
MAP, $\alpha=2$ & $(N_j+1)/(N+K)$ & igual que la fila anterior \\
\bottomrule
\end{tabular}
\end{table}

\addtocounter{definicion}{1}%


% =====================================================================
\chapter{Teoremas límite y estimación por Monte Carlo}
\label{tema:5}
% =====================================================================

Ni la pérdida de un LLM ni su exactitud en un banco de pruebas pueden calcularse exactamente: son esperanzas respecto de una distribución desconocida y se estiman con promedios. Este tema demuestra que esos promedios convergen (ley de los grandes números), cuantifica su error (desigualdad de Chebyshev y teorema central del límite) y lo aplica al método de Monte Carlo, a los intervalos de confianza y a los minilotes del entrenamiento. Se sigue el capítulo~5 de~\cite{bt}, las lecciones~19 y~20 de~\cite{mit6041}, la clase~6b de~\cite{mit1805} y el capítulo~10 de~\cite{bh}; la §2.6 de~\cite{d2l} trata el mismo material en el contexto del aprendizaje profundo.

\section{Desigualdades de Markov y Chebyshev}

\begin{teorema}[Desigualdad de Markov]\label{teo:markov}
Si $Y\ge 0$ y $a>0$, entonces $\PP(Y\ge a)\le\E[Y]/a$.
\end{teorema}

\begin{teorema}[Desigualdad de Chebyshev]\label{teo:chebyshev}
Si $X$ tiene media $\mu$ y varianza $\sigma^{2}$, para todo $\varepsilon>0$,
\[
\PP\bigl(\abs{X-\mu}\ge\varepsilon\bigr)\le\frac{\sigma^{2}}{\varepsilon^{2}}.
\]
\end{teorema}

Las dos desigualdades valen para cualquier distribución con los momentos indicados; a cambio, suelen ser pesimistas~\cite[§5.1]{bt}.

\section{Ley débil de los grandes números}

\begin{definicion}[Convergencia en probabilidad]
Una sucesión $Y_1,Y_2,\dots$ \emph{converge en probabilidad} a la constante $c$ si, para todo $\varepsilon>0$, $\lim_{n\to\infty}\PP\bigl(\abs{Y_n-c}\ge\varepsilon\bigr)=0$.
\end{definicion}

\begin{teorema}[Ley débil de los grandes números]\label{teo:lgn}
Sean $X_1,X_2,\dots$ \iid\ con media $\mu$ y varianza finita $\sigma^{2}$. Entonces, para todo $\varepsilon>0$,
\[
\PP\bigl(\abs{\Xbar_n-\mu}\ge\varepsilon\bigr)\le\frac{\sigma^{2}}{n\varepsilon^{2}}\xrightarrow[n\to\infty]{}0,
\]
es decir, $\Xbar_n$ converge en probabilidad a $\mu$.
\end{teorema}

La ley fuerte de los grandes números afirma que, de hecho, $\Xbar_n\to\mu$ con probabilidad~1; su demostración excede este módulo~\cite[§5.5]{bt}.

\section{El método de Monte Carlo}

\begin{definicion}[Estimador de Monte Carlo]\label{def:montecarlo}
Para aproximar $I=\E[g(X)]$, se generan $X_1,\dots,X_n$ \iid\ con la distribución de $X$ y se toma $\hat I_n=\frac1n\sum_{i=1}^{n}g(X_i)$. Por el \cref{cor:media}, $\hat I_n$ es insesgado y su \emph{error típico} es $\sigma_g/\sqrt n$, donde $\sigma_g^{2}=\Var(g(X))$.
\end{definicion}

\addtocounter{definicion}{1}%

La ley $1/\sqrt n$ no depende de la dimensión del problema, lo que hace a Monte Carlo competitivo para integrales en muchas dimensiones, pero sí la hace lenta: el error se reduce a la mitad solo cuadruplicando el número de muestras.

\addtocounter{definicion}{1}%

\section{Teorema central del límite}

La ley de los grandes números dice que $\Xbar_n-\mu\to 0$; el teorema central del límite describe la forma de la distribución de ese error cuando se amplía por el factor $\sqrt n$.

\begin{teorema}[Teorema central del límite]\label{teo:tcl}
Sean $X_1,X_2,\dots$ \iid\ con media $\mu$ y varianza $\sigma^{2}\in(0,\infty)$, y sea $S_n=X_1+\dots+X_n$. Entonces, para todo $z\in\R$,
\[
\lim_{n\to\infty}\PP\Bigl(\frac{S_n-n\mu}{\sigma\sqrt n}\le z\Bigr)=\Phi(z).
\]
\end{teorema}

La demostración habitual usa funciones generatrices de momentos o funciones características~\cite[§5.4]{bt}. Lo notable es que no se exige que los $X_i$ sean normales: la normalidad es la conclusión, no la hipótesis. Así se explica la ubicuidad de la normal: el ruido térmico, los errores de medida o la suma de muchos pesos pequeños son sumas de muchos efectos independientes.

\begin{corolario}[Aproximación normal de la media]\label{cor:aprox}
Para $n$ grande, $\Xbar_n$ es aproximadamente $\N(\mu,\sigma^{2}/n)$ y, en particular,
\[
\PP\Bigl(\abs{\Xbar_n-\mu}\ge z\,\frac{\sigma}{\sqrt n}\Bigr)\approx 2\bigl(1-\Phi(z)\bigr).
\]
\end{corolario}

\begin{table}[h]
\centering\small
\caption{Tres herramientas aplicadas al \cref{ej:pi} con $n=10^{4}$ y un error de $0{,}05$, que equivale a $3{,}04$ errores típicos.}\label{tab:herramientas}
\begin{tabular}{@{}lll@{}}
\toprule
Herramienta & Afirmación & Carácter \\
\midrule
Ley de los grandes números & $\hat\pi_n\to\pi$ & cualitativa \\
Chebyshev & $\PP(\abs{\hat\pi_n-\pi}\ge 0{,}05)\le 0{,}108$ & cota rigurosa, pesimista \\
Teorema central del límite & $\PP(\abs{\hat\pi_n-\pi}\ge 0{,}05)\approx 0{,}0023$ & aproximación realista \\
\bottomrule
\end{tabular}
\end{table}

\section{Intervalos de confianza y comparación de modelos}

\begin{definicion}[Intervalo de confianza aproximado]
Sea $\hat p$ la proporción de aciertos de un modelo en $n$ preguntas independientes. El intervalo
\[
\hat p\pm 1{,}96\sqrt{\frac{\hat p(1-\hat p)}{n}}
\]
es un \emph{intervalo de confianza aproximado del $95\,\%$} para la exactitud verdadera $p$: por el \cref{cor:aprox}, contiene a $p$ con probabilidad aproximada $0{,}95$~\cite[clase~22]{mit1805}.
\end{definicion}

\addtocounter{definicion}{1}%

\section{Minilotes y gradiente estocástico}

La función de pérdida de entrenamiento es una esperanza, $L(\theta)=\E[\lambda(\theta;X)]$, donde $\lambda(\theta;x)$ es la pérdida en el ejemplo $x$ y $X$ sigue la distribución de los datos. Un \emph{minilote} es una muestra $X_1,\dots,X_B$ extraída al azar, y la pérdida del minilote es
\[
\hat L_B(\theta)=\frac1B\sum_{i=1}^{B}\lambda(\theta;X_i).
\]

\begin{proposicion}\label{prop:minilote}
Si los ejemplos del minilote son \iid\ y $\sigma^{2}=\Var(\lambda(\theta;X))$, entonces $\E[\hat L_B(\theta)]=L(\theta)$ y la desviación típica de $\hat L_B(\theta)$ es $\sigma/\sqrt B$. Si la derivada y la esperanza pueden intercambiarse, $\nabla_\theta\hat L_B$ es también un estimador insesgado de $\nabla_\theta L$, cuya covarianza decrece como $1/B$.
\end{proposicion}

\addtocounter{definicion}{1}%

El descenso de gradiente estocástico (módulo~7) avanza con estos gradientes ruidosos. Por el \cref{teo:tcl}, el ruido es aproximadamente normal, y el tamaño de lote es un compromiso entre la calidad de cada paso y su coste.



\backmatter
% =====================================================================
\begin{thebibliography}{99}
\addcontentsline{toc}{chapter}{Bibliografía}
\markboth{Bibliografía}{}
\small

\bibitem{bt}
D.~P. Bertsekas y J.~N. Tsitsiklis.
\newblock \emph{Introduction to Probability}, 2.ª ed.
\newblock Athena Scientific, 2008.
\newblock Texto del curso 6.041 del MIT.

\bibitem{bishop}
C.~M. Bishop.
\newblock \emph{Pattern Recognition and Machine Learning}.
\newblock Springer, 2006.

\bibitem{bh}
J.~K. Blitzstein y J.~Hwang.
\newblock \emph{Introduction to Probability}, 2.ª ed.
\newblock CRC Press, 2019.
\newblock Texto del curso Stat~110 de Harvard.

\bibitem{stat110}
J.~K. Blitzstein.
\newblock Stat 110: Probability.
\newblock Harvard University. \emph{Strategic Practice and Homework}, otoño de 2011.
\newblock \url{https://stat110.hsites.harvard.edu/strategic-practice-problems}

\bibitem{mml}
M.~P. Deisenroth, A.~A. Faisal y C.~S. Ong.
\newblock \emph{Mathematics for Machine Learning}.
\newblock Cambridge University Press, 2020.
\newblock \url{https://mml-book.github.io/}

\bibitem{glorot}
X.~Glorot e Y.~Bengio.
\newblock Understanding the difficulty of training deep feedforward neural networks.
\newblock En \emph{Proceedings of the Thirteenth International Conference on Artificial Intelligence and Statistics (AISTATS)}, PMLR~9:249--256, 2010.
\newblock \url{https://proceedings.mlr.press/v9/glorot10a.html}

\bibitem{he}
K.~He, X.~Zhang, S.~Ren y J.~Sun.
\newblock Delving deep into rectifiers: surpassing human-level performance on ImageNet classification.
\newblock En \emph{Proceedings of the IEEE International Conference on Computer Vision (ICCV)}, 2015.
\newblock \url{https://arxiv.org/abs/1502.01852}

\bibitem{mit1805}
J.~Orloff y J.~Bloom.
\newblock 18.05: Introduction to Probability and Statistics.
\newblock Massachusetts Institute of Technology, MIT OpenCourseWare, primavera de 2022.
\newblock \url{https://ocw.mit.edu/courses/18-05-introduction-to-probability-and-statistics-spring-2022/}

\bibitem{mit6041}
J.~N. Tsitsiklis.
\newblock 6.041SC: Probabilistic Systems Analysis and Applied Probability.
\newblock Massachusetts Institute of Technology, MIT OpenCourseWare, otoño de 2013.
\newblock \url{https://ocw.mit.edu/courses/6-041sc-probabilistic-systems-analysis-and-applied-probability-fall-2013/}

\bibitem{mit6041f10}
J.~N. Tsitsiklis.
\newblock 6.041: Probabilistic Systems Analysis and Applied Probability.
\newblock Massachusetts Institute of Technology, MIT OpenCourseWare, otoño de 2010. \emph{Problem Set}~1.
\newblock \url{https://ocw.mit.edu/courses/6-041-probabilistic-systems-analysis-and-applied-probability-fall-2010/}

\bibitem{xie}
S.~M. Xie, A.~Raghunathan, P.~Liang y T.~Ma.
\newblock An explanation of in-context learning as implicit Bayesian inference.
\newblock En \emph{International Conference on Learning Representations (ICLR)}, 2022.
\newblock \url{https://arxiv.org/abs/2111.02080}

\bibitem{d2l}
A.~Zhang, Z.~C. Lipton, M.~Li y A.~J. Smola.
\newblock \emph{Dive into Deep Learning}, versión 1.0.3.
\newblock Cambridge University Press, 2023.
\newblock \url{https://d2l.ai/}

\end{thebibliography}


\end{document}
